\documentclass[../../main.tex]{subfiles}% 注意这里的文件路径不能用 ./main.tex ,否则用latexmk编译子文件会报错
\graphicspath{{\subfix{./image/}}} % 指定图片目录,后续可以直接使用图片文件名
% 注意这里的文件路径不能用 ../../image/ ,否则用latexmk编译子文件会报错

% 例如:
% \begin{figure}[H]
% \centering
% \includegraphics[scale=0.3]{图.png}
% \caption{}
% \label{figure:图}
% \end{figure}
% 注意:上述\label{}一定要放在\caption{}之后,否则引用图片序号会只会显示??.

\begin{document}

\section{平面曲线}

平面曲线是挠率$\tau\equiv0$的空间曲线，具有独立的理论体系。

设平面曲线$\bm{r}(s)=(x(s),y(s))$，$s$为弧长，单位切向量$\bm{\alpha}(s)=(x'(s),y'(s))$，将$\bm{\alpha}(s)$逆时针旋转$90^\circ$得**法向量**：
\[
\bm{\beta}(s) = \big(-y'(s),x'(s)\big).
\]


\begin{definition}[相对曲率]
平面曲线的**相对曲率**定义为
\[
\kappa_r(s) = \bm{\alpha}'(s)\cdot\bm{\beta}(s),
\]
满足$\bm{\alpha}'(s)=\kappa_r(s)\bm{\beta}(s)$。
\end{definition}

一般参数$t$下的相对曲率公式：
\[
\kappa_r(t) = \dfrac{x'(t)y''(t)-x''(t)y'(t)}{\big(x'(t)^2+y'(t)^2\big)^{\frac{3}{2}}}.
\]


\begin{definition}[旋转指标]
光滑简单闭曲线的单位切向量绕曲线一周的旋转角除以$2\pi$，称为**旋转指标**，记为$i(C)$。
\end{definition}

\begin{theorem}[旋转指标定理]
平面上连续可微的简单闭曲线的旋转指标$i(C)=\pm1$（正向为$+1$，负向为$-1$）。
\end{theorem}

\begin{example}
圆、椭圆的旋转指标为$1$，简单闭曲线的旋转指标恒为$\pm1$。
\end{example}

\begin{exercise}
\begin{enumerate}
\item 求摆线$\bm{r}(t)=\big(a(t-\sin t),a(1-\cos t)\big)$的相对曲率。
\item 证明：平面简单闭曲线的旋转指标为$\pm1$。
\end{enumerate}
\end{exercise}




\end{document}